\documentclass[UTF8,fontset=none,11pt,a4paper]{ctexart}
\usepackage{ipara-handbook}
\handbooksetup{NET-07 拥塞控制}{408 · 计算机网络 NET-07}{Feedback · cwnd · AIMD · Reno}
\providecommand{\tightlist}{\setlength{\itemsep}{0pt}\setlength{\parskip}{0pt}}
\providecommand{\passthrough}[1]{#1}
\begin{document}
\handbooktitle{拥塞控制}{在不知道路径容量时闭环试探}{408 · NET-07}
\section{0. 本册定位}\label{0-ux672cux518cux5b9aux4f4d}

本 Topic
回答：大量独立发送者共享有限链路与队列、又看不到全局负载时，端点怎样用反馈调节
offered load，避免网络进入高排队、高丢包、低有效吞吐的失控状态？

本册拥有 congestion
signal、\passthrough{\lstinline!cwnd!}、\passthrough{\lstinline!ssthresh!}、slow
start、congestion avoidance、AIMD、fast retransmit/recovery 的 408
经典模型，以及效率、公平、时延和稳定性的权衡。

本册使用 TCP sequence/ACK/connection
state，见\href{../06_传输层_端点_UDP与TCP状态机/README.md}{传输层}；可靠性与接收端流控分别不等同于
congestion control。

\section{1.
根本问题：总输入可以超过共享服务能力}\label{1-ux6839ux672cux95eeux9898ux603bux8f93ux5165ux53efux4ee5ux8d85ux8fc7ux5171ux4eabux670dux52a1ux80fdux529b}

设某瓶颈链路服务率为 \(C\)，多个流的总到达速率为 \(\lambda\)：

\begin{itemize}
\tightlist
\item
  \(\lambda<C\)：队列通常能被清空；
\item
  \(\lambda\approx C\)：小波动就会积累排队；
\item
  \(\lambda>C\) 持续存在：队列增长，最终丢包或时延失控。
\end{itemize}

朴素方案是每个 sender 都尽可能快地发送。单个 sender
看起来在提高自身吞吐，但所有 sender 的独立最优叠加会破坏共享资源：

\[
\boxed{
\text{More offered load}
\to \text{Queue growth}
\to \text{Delay/loss}
\to \text{Retransmission}
\to \text{Even more load}
}
\]

拥塞控制必须打断这个正反馈。

\section{2. 拥塞是路径状态，不是单个 packet
的属性}\label{2-ux62e5ux585eux662fux8defux5f84ux72b6ux6001ux4e0dux662fux5355ux4e2a-packet-ux7684ux5c5eux6027}

发送端通常不知道：瓶颈在哪里、可用容量多少、同时有多少竞争流。它只能观察反馈：

\begin{itemize}
\tightlist
\item
  loss/timeout；
\item
  duplicate ACK pattern；
\item
  increasing RTT/queue delay；
\item
  ECN 等显式标记；
\item
  ACK 到达节奏。
\end{itemize}

然后更新本地发送约束：

\[
\boxed{
\text{Probe}
\to \text{Observe feedback}
\to \text{Infer path state}
\to \text{Adjust cwnd}
\to \text{Probe again}
}
\]

这是反馈控制，不是一次算出带宽后永久使用。

\section{\texorpdfstring{3. \texttt{cwnd}
的对象与单位}{3. cwnd 的对象与单位}}\label{3-cwnd-ux7684ux5bf9ux8c61ux4e0eux5355ux4f4d}

\passthrough{\lstinline!cwnd!} 是 sender 为该 connection 维护的
congestion window，限制网络中允许存在的未确认数据量。TCP 实际发送还受
receiver window：

\[
FlightSize\le \min(cwnd,rwnd).
\]

408 题目有时用 MSS 个数表示窗口，有时用 byte。计算前必须写单位；不能把
\(\text{cwnd} = 8\) 自动解释成 8 byte 或 8 segment。

若 sender 希望充分利用路径，理想在途量与路径 BDP
同量级；大幅超过时，多余数据主要进入 queue，而不是让传播管道变长。

\section{4. Slow
Start：在未知容量下快速探测}\label{4-slow-startux5728ux672aux77e5ux5bb9ux91cfux4e0bux5febux901fux63a2ux6d4b}

连接开始或严重拥塞后，若从 1 MSS 线性增加，找到高速路径容量太慢。Slow
Start 借 ACK clock 让 \passthrough{\lstinline!cwnd!} 在每个 RTT 约翻倍：

\begin{lstlisting}
cwnd = 1 MSS
one RTT of ACKs -> about 2 MSS
next RTT        -> about 4 MSS
next RTT        -> about 8 MSS
...
\end{lstlisting}

名称中的 ``slow'' 是相对一开始直接发送巨大 burst；它的增长在 RTT
维度上是指数级。

当 \passthrough{\lstinline!cwnd!} 到达
\passthrough{\lstinline!ssthresh!} 附近，切换到 congestion
avoidance。\passthrough{\lstinline!ssthresh!}
不是路径真实容量，而是根据历史拥塞更新的本地分界估计。

\section{5. Congestion
Avoidance：加性试探}\label{5-congestion-avoidanceux52a0ux6027ux8bd5ux63a2}

接近已知可用范围后，继续指数增长会频繁过冲。经典 congestion avoidance 让
\passthrough{\lstinline!cwnd!} 每 RTT 大约增加 1 MSS：

\[
\Delta cwnd\approx 1\ MSS/RTT.
\]

每个 ACK 的具体增量可按当前 \passthrough{\lstinline!cwnd!}
分摊，使一整个窗口被确认后累计约 1
MSS。核心不是背某个实现公式，而是识别增长从 multiplicative 变为
additive。

\section{6. Congestion event
后为什么要乘性下降}\label{6-congestion-event-ux540eux4e3aux4ec0ux4e48ux8981ux4e58ux6027ux4e0bux964d}

若检测到拥塞仍只减一个固定小量，多个 flow
可能长时间把总负载维持在瓶颈之上。经典 AIMD：

\[
\boxed{\text{Additive Increase} + \text{Multiplicative Decrease}}
\]

用缓慢线性增长持续探测剩余容量，用比例下降快速释放大量在途负载。AIMD
在理想同质条件下还具有向效率与公平状态收敛的几何直觉，但真实
RTT、MSS、应用模式和算法差异会影响公平性。

\section{7. Timeout 与 duplicate ACK
表示不同强度的证据}\label{7-timeout-ux4e0e-duplicate-ack-ux8868ux793aux4e0dux540cux5f3aux5ea6ux7684ux8bc1ux636e}

\subsection{7.1
Timeout：反馈闭环严重中断}\label{71-timeoutux53cdux9988ux95edux73afux4e25ux91cdux4e2dux65ad}

RTO 到期意味着对应数据在足够长时间内没有获得预期确认。经典 408/Reno
模型把它视为较严重拥塞：更新 \passthrough{\lstinline!ssthresh!}，将
\passthrough{\lstinline!cwnd!} 降到较小初值，重新 slow start。

\subsection{7.2 Three duplicate
ACKs：后续数据仍在通过}\label{72-three-duplicate-acksux540eux7eedux6570ux636eux4ecdux5728ux901aux8fc7}

连续 duplicate ACK 表明接收端反复看见同一个缺口，但缺口之后的一些
segments 已到达。这通常比 timeout
提供更细的信息：路径仍在交付数据，只是某个 segment 可能丢失。

经典模型由此产生：

\begin{itemize}
\tightlist
\item
  fast retransmit：不等 RTO，重传推定丢失 segment；
\item
  fast recovery：不把发送速率退回最初状态，而按相对温和方式降低并继续。
\end{itemize}

具体 \(\text{cwnd}/\text{ssthresh}\) 数值更新必须服从题目指定的
Tahoe/Reno 教材模型；现代 TCP 实现有多种 congestion-control
algorithms，不能把一张 408 状态图冒充所有工程实现。

\section{8. ACK
clock：反馈不仅说``收到''，还提供节奏}\label{8-ack-clockux53cdux9988ux4e0dux4ec5ux8bf4ux6536ux5230ux8fd8ux63d0ux4f9bux8282ux594f}

当 ACK 随数据穿过瓶颈返回，ACK 到达节奏近似反映路径的交付节奏。sender 按
ACK 释放新数据，可以避免无节制 burst。

但 ACK compression、delayed ACK、reverse-path congestion
等会扭曲时钟。ACK clock 是有用反馈机制，不是精确的全局容量测量仪。

\section{9. Queue
的双重身份}\label{9-queue-ux7684ux53ccux91cdux8eabux4efd}

Queue 不是纯粹坏事：短队列吸收
burst，使链路持续忙碌。问题是持久过量队列：

\begin{itemize}
\tightlist
\item
  latency 增大；
\item
  timeout 和交互性能恶化；
\item
  buffer overflow 导致 loss；
\item
  新流和短流被已有大队列拖累。
\end{itemize}

因此优化目标不是``绝不排队''或``永不丢包''，而是在
utilization、delay、loss、fairness 和 stability 之间取舍。

\section{10. Network-assisted 与 end-to-end
feedback}\label{10-network-assisted-ux4e0e-end-to-end-feedback}

拥塞控制先分“事前约束”与“事后反馈”。事前/open-loop 策略在观察本次拥塞前限制系统行为，例如接纳控制、资源预留、流量整形、调度与丢弃策略；反馈/closed-loop 策略检测 queue、delay、loss 或显式标记，再把信号送回能降低 offered load 的主体。前者降低进入危险区的概率，后者适应实际运行状态；二者可以同时存在。

再按反馈 Owner 分层：router/network layer 可调度队列、丢弃或标记 packet，也可在特定体系中反馈显式速率；transport endpoint 根据 ACK/loss/ECN/RTT 等调整发送窗口或速率。网络层“发现/标记拥塞”不自动完成端点减速，端点算法也不能直接观察每台 router 的完整队列真相。

\begin{itemize}
\tightlist
\item
  loss-based control：从丢包/ACK 推断拥塞，部署简单但信号较晚；
\item
  ECN：网络设备在丢包前标记拥塞，端点仍负责降低发送；
\item
  delay-based control：从 RTT/queue delay
  变化提前推断，但需要区分路径基线和噪声；
\item
  explicit rate
  等更强网络协助模型直接给出速率或资源信息，但需要基础设施支持。
\end{itemize}

408 核心以经典 TCP loss-based
模型为主；扩展用于解释反馈设计空间，不替代考试指定算法。

因此网络层拥塞与 TCP 拥塞控制不是两种互斥现象：前者描述共享 router/link 的资源状态与网络可提供的控制动作，后者是端点借助有限反馈调节 offered load 的一种闭环实现。

\section{11. 408 状态推进：把题设算法写成转移函数}

先固定单位（byte 或 MSS）、初值、`ssthresh`、ACK 粒度与题设版本。然后对每个事件应用：

\begin{handbooktable}{L{2.6cm} Y Y Y}
\toprule
事件 & 进入前状态 & 经典教材动作 & 停止/下一分支\\
\midrule
new ACK，$cwnd<ssthresh$ & slow start & 每 ACK 增长；一 RTT 近似翻倍 & 达阈值转 avoidance\\
new ACK，$cwnd\ge ssthresh$ & avoidance & 一 RTT 近似增加 1 MSS & 继续探测\\
timeout & 任一发送阶段 & $ssthresh\leftarrow$ 题设比例；$cwnd\leftarrow$ 小初值 & 重传并回 slow start\\
three duplicate ACKs，Tahoe & 检测单缺口 & fast retransmit；窗口降到小初值 & slow start\\
three duplicate ACKs，Reno & 检测单缺口 & fast retransmit；进入/维持 fast recovery & new ACK 后回 avoidance\\
\bottomrule
\end{handbooktable}

具体加一、减半、初始窗口和 recovery 膨胀数值以题目定义为准。手册拥有的是“状态、事件、反馈、转移”模型，不把某一版本常数写成所有 TCP 的永恒事实。

\begin{examplebox}{轮次与逐 ACK 不能混算}
若题目说 slow start 每收到一个确认新数据的 ACK 增加 1 MSS，则一个窗口内的多个 ACK 累计造成下一 RTT 可发送量近似翻倍；不能把“每 ACK 加 1”误写成“每 ACK 翻倍”。题目若直接按 RTT 给窗口序列，则不要再重复逐 ACK 累加。
\end{examplebox}

\section{12. 概念边界}\label{11-ux6982ux5ff5ux8fb9ux754c}

\begin{longtable}[]{@{}lclll@{}}
\toprule\noalign{}
概念 A & ≠ & 概念 B & 真正区别与题目信号 & 混淆后果 \\
\midrule\noalign{}
\endhead
\bottomrule\noalign{}
\endlastfoot
Congestion & ≠ & Receiver overflow & 路径共享资源过载；接收应用/缓冲不足
& \passthrough{\lstinline!cwnd!} 与 \passthrough{\lstinline!rwnd!}
更新写反 \\
\passthrough{\lstinline!cwnd!} & ≠ & \passthrough{\lstinline!rwnd!} &
sender 的路径推断；receiver 的容量通告 & 错判谁持有和谁反馈 \\
\passthrough{\lstinline!ssthresh!} & ≠ & Bottleneck capacity &
本地算法阶段阈值；路径容量是外部动态状态 & 把阈值当精确带宽 \\
Timeout & ≠ & Certain congestion & timeout
是强推断，也可能受非拥塞丢失/延迟影响 & 把观测当全局事实 \\
Duplicate ACK & ≠ & New data ACK & 前者重复声明同一连续前缀缺口 & ACK
数与窗口推进错误 \\
Flow rate & ≠ & Goodput & 前者可含重传/首部；后者只数应用有效数据 &
拥塞时``发送更多''误判为性能更好 \\
Slow start & ≠ & Linear slow growth & RTT 维度近似指数增长 &
窗口轮次推演错误 \\
Classic Reno model & ≠ & All modern TCP & 408 模型是一个算法实例 &
用旧状态图断言现实协议唯一行为 \\
\end{longtable}

\section{12.
做题调用协议}\label{12-ux505aux9898ux8c03ux7528ux534fux8bae}

\begin{enumerate}
\def\labelenumi{\arabic{enumi}.}
\tightlist
\item
  标出
  \passthrough{\lstinline!cwnd!}、\passthrough{\lstinline!ssthresh!}、\passthrough{\lstinline!rwnd!}
  的单位和初值；
\item
  画 RTT 轮次或逐 ACK 事件，按题目粒度更新；
\item
  判断当前处于 slow start、avoidance 还是 recovery；
\item
  区分 timeout 与 duplicate ACK 触发的教材分支；
\item
  计算实际发送上限时取 \(\min(\text{cwnd}, \text{rwnd})\) 并扣除
  in-flight；
\item
  明确题目采用 Tahoe、Reno 或自定义规则；
\item
  用 BDP、queue 和 goodput 检查窗口变化是否有物理意义。
\end{enumerate}

\section{13. 贯穿母例：为什么 1、2、4、8
后不能永远翻倍}\label{13-ux8d2fux7a7fux6bcdux4f8bux4e3aux4ec0ux4e48-1248-ux540eux4e0dux80fdux6c38ux8fdcux7ffbux500d}

假设瓶颈在途与队列可暂时容纳约 10 MSS。sender 从 1 MSS slow start：

\begin{lstlisting}
1 -> 2 -> 4 -> 8: path may absorb
8 -> 16: offered flight overshoots available capacity
-> queue grows / loss signal appears
-> sender reduces cwnd and records ssthresh
-> later probes additively near the observed boundary
\end{lstlisting}

Slow start 解决``完全不知道容量时怎样快速接近''；AIMD
解决``接近边界后怎样持续试探并在过冲时退让''。二者不是两套互不相关的背诵规则，而是探索阶段不同。

\section{14. 高频 First
Divergence}\label{14-ux9ad8ux9891-first-divergence}

\begin{itemize}
\tightlist
\item
  把 \passthrough{\lstinline!cwnd!} 当 receiver 通告值：状态 owner 错；
\item
  每收到一个 ACK 就让 \passthrough{\lstinline!cwnd!} 翻倍：混淆逐 ACK
  增量与逐 RTT 效果；
\item
  到 \passthrough{\lstinline!ssthresh!} 后停止增长：忘记 avoidance
  仍在探测；
\item
  timeout 与三次重复 ACK 使用同一分支：丢失证据强度分层；
\item
  只算发送速率不算 goodput：重传负载被当成有效数据；
\item
  未读题设就套 Reno：没有确认教材模型。
\end{itemize}

\section{15.
一页压缩与复原问题}\label{15-ux4e00ux9875ux538bux7f29ux4e0eux590dux539fux95eeux9898}

\[
\boxed{
\text{Probe capacity}
\to \text{Read ACK/loss/ECN/delay}
\to \text{Infer congestion}
\to \text{Adjust cwnd}
\to \text{Probe again}
}
\]

\begin{enumerate}
\def\labelenumi{\arabic{enumi}.}
\tightlist
\item
  为什么 congestion control 必须是闭环而不是固定速率？
\item
  Slow Start 为什么既``慢''又近似指数增长？
\item
  AIMD 的加性与乘性分别解决什么？
\item
  duplicate ACK 为什么通常比 timeout 提供更细的路径信息？
\item
  Queue 在什么范围内有益，何时开始伤害系统？
\end{enumerate}

\section{16.
来源与校正说明}\label{16-ux6765ux6e90ux4e0eux6821ux6b63ux8bf4ux660e}

\begin{itemize}
\tightlist
\item
  归档材料没有独立拥塞控制正文，本册不是从术语表复制而来；
\item
  \passthrough{\lstinline!cwnd!}、\passthrough{\lstinline!rwnd!}
  共同限制在途数据，以及 slow start/congestion avoidance/fast
  retransmit/recovery 的经典边界依据
  \href{https://www.rfc-editor.org/rfc/rfc5681.html}{RFC 5681} 校正；
\item
  具体窗口初值和现代算法会演进，本册将工程现实与 408
  指定模型分层，不把版本相关常数写成永恒机制。
\end{itemize}
\end{document}
